\chapter{Conclusiones}

En esta tesina se planteó el problema de detectar automáticamente inconsistencias en las firmas manuscritas de los telegramas electorales utilizados en el escrutinio provisorio con asistencia digital en la Provincia de Santa Fe. A diferencia del enfoque clásico de VAFM, el objetivo no fue identificar falsificaciones intencionales, sino errores de confección (ausencia de firmas o firmas no coincidentes entre páginas) que derivan en retransmisiones y trabajo adicional en el centro de cómputo. El problema se formalizó como una tarea de clasificación binaria sobre pares de imágenes de firmas, definiendo cuidadosamente las métricas de evaluación (\textit{precision}, \textit{recall}, curvas ROC / Precision–Recall) y el rol asimétrico de los errores según el uso práctico del sistema. 

Desde el punto de vista metodológico, se construyó primero un conjunto de datos propio a partir de telegramas reales, aprovechando que el diseño de los formularios es conocido de antemano. Se diseñó un pipeline donde, a partir de fotografías de telegramas, se recortan automáticamente las zonas de firmas y se entrena un modelo intermedio para detectar presencia/ausencia de firma, que permitió limpiar y estructurar el conjunto de datos definitivo. Sobre este \textit{dataset} se compararon dos familias de enfoques: (i) métodos “clásicos” basados en extracción manual de características (grillas y descriptores inspirados en \cite{tesis-parodi}) combinados con clasificadores como Random Forest y SVM; y (ii) modelos de Aprendizaje Profundo basados en redes convolucionales y arquitecturas siamesas para aprendizaje de representaciones y métricas de similitud. 

Los experimentos muestran que las redes siamesas convolucionales, entrenadas con funciones de costo específicas para aprendizaje de métricas (como \textit{triplet loss} en variantes \textit{online} y \textit{offline}), logran resultados claramente superiores a los enfoques basados en características manuales, tanto en conjuntos de datos estándar de firmas como en el escenario particular de los telegramas. Se alcanzan configuraciones de operación donde el modelo mantiene un \textit{recall} en torno a 0.8–0.9 con alta \textit{precision}, lo que en términos prácticos se traduce en una reducción significativa del número de retransmisiones por errores de firmas, manteniendo a la vez un volumen limitado de falsas alarmas aceptable para el flujo de trabajo del centro de cómputo. 

Como principales limitaciones se destaca la dependencia del modelo respecto del formato de los telegramas (posición fija de las firmas y calidad de las imágenes), la ausencia de falsificaciones intencionales en el conjunto de datos (el foco está puesto en errores de confección) y la necesidad de contar con volúmenes suficientes de datos etiquetados para explorar arquitecturas más grandes. En el capítulo de Trabajo futuro se discuten posibles extensiones, entre ellas: el uso de representaciones aprendidas mediante modelos generativos (autoencoders o GANs), la exploración de arquitecturas siamesas y de \textit{metric learning} más modernas, y la adaptación del pipeline a otros problemas de similitud de imágenes, como verificación de identidad en documentos o análisis de otras marcas manuscritas. Estas líneas abren el camino para seguir mejorando la robustez y capacidad de generalización del sistema propuesto.


